\section{Related Work}
\vpara{Evaluation of LLMs.}
The general capabilities of self-supervised~\citep{liu2021self} LLMs~\citep{GPT3,chowdhery2022palm,zhang2022opt,scao2022bloom,zeng2022glm,touvron2023llama}, especially those chat-aligned ones~\citep{ouyang2022training,Claude,openai2023gpt}, have refreshed people's impression on deep learning systems and significantly transcended the conventional scope of NLP evaluation.
It thus makes the evaluation of LLMs an urgent and challenging problem.
Compared to previous efforts focusing on a subset of specified tasks~\citep{wang2019superglue,wangglue,gehrmann2021gem}, an increasing number of benchmarks are including broader spectra of tasks and datasets~\citep{hendrycksmeasuring,liang2022holistic,srivastava2023beyond} in the evaluation.
However, most of them are still limited to traditional tasks and thus fail to evaluate LLMs' open-ended generation, multi-round interaction, and ability to act as agents.

\vpara{LLM-as-Agent.}
In pre-LLM era, text game environments such as TextWorld~\citep{cote2019textworld}, Jericho~\citep{hausknecht2020interactive}, and LIGHT~\citep{urbanek2019learning} are dominant in language agent study which bases on BERT~\citep{devlin2019bert} and reinforcement learning.
With the advent of LLMs, the study of LLM agents begins to thrive~\citep{huang2022language}, especially after Chain-of-Thought~\citep{wei2022chain} came out.
ReAct~\citep{yao2023react} is a pioneer work to combine CoT reasoning and actions in agent tasks.
Later, a multitude of advanced reasoning strategies~\citep{kim2023language,shinn2023reflexion,wang2023describe,liu2023llm+,yao2023tree,gu-etal-2023-dont} and applications including frameworks~\citep{richards2023auto,nakajima2023babyagi,agentgpt} and multi-agents~\citep{Park2023GenerativeAI,Hong2023MetaGPTMP,Wu2023AutoGenEN} for LLM-as-Agent have emerged and arouse much public interest.
Nevertheless, limited datasets and models and available on the topic, without a standard and comprehensive benchmark.
\model presents the first systematic benchmark for evaluating LLM-as-Agent with a broad coverage of tasks and available LLMs.
Additionally, it also initiates the idea of adopting agent tasks to measure LLM performance.

\vpara{Evaluating LLMs in Executive Environments.}
As LLMs become increasingly capable of real-world challenges, there is also a trend to evaluate them in executive environments rather than static datasets.
Besides text games (e.g., ALFWorld~\citep{shridhar2020alfworld}), another main stream of works lies in code execution.
APPS~\citep{hendrycks2021measuring}, HumanEval~\citep{chen2021evaluating} and MBPP~\citep{austin2021program} pioneer the effort to evaluate code LLMs for functional correctness instead of text similarity.
The paradigm has been later widely recognized and adopted in following works~\citep{li2022competition,zheng2023codegeex,nijkamp2023codegen}.
However, few previous code evaluation frameworks consider multi-round interactions.
A concurrent work InterCode~\citep{yang2023intercode} releases a framework that allows evaluation of interaction between models and Bash and SQL environments, which are similar to OS and DB tasks in \model. 

\hhide{
However, it is different from \model in several aspects: 1) It displays no command and environment output histories to agent and thus limits agents' capability. 2) Only model's execution output and the difference of file system are checked, leaving other systems (e.g., network system, user system) unverfiable. 3) Each round an independent shell is launched, disallowing advanced commands such as variable declaration.

Most importantly, prior efforts on the topic do not present such a diverse array of task environments beyond coding.
Few previous works would also have comparably extensive coverage of LLM testing as \model does.
}